% chapters/13/section-02.tex
\section{Collider dan Rigidbody}
\label{sec:collider-rigidbody}

Dua komponen inilah yang menentukan fisika sebuah level. Bagian ini menjelaskan
pilihan yang tersedia pada keduanya, sesuai indikator \ind{b}.

\subsection{Bentuk Batas Benda}

Komponen \istilah{Collider} menentukan bentuk batas benda untuk keperluan
tabrakan. Bentuk yang tersedia bergantung pada pandangan yang dipakai, dan
yang paling sering dipakai pada permainan dua dimensi adalah
\istilah{Box Collider 2D}, \istilah{Circle Collider 2D}, dan
\istilah{Capsule Collider 2D}.

Tiga hal perlu diperhatikan pada pemakaiannya.

\begin{enumerate}
  \item \textbf{Bentuk batas tidak harus sama dengan gambarnya.} Sebuah
        karakter bergambar rumit cukup diberi satu \istilah{Capsule Collider}
        yang melingkupinya. Cara itu jauh lebih murah, dan hasilnya sering kali
        justru lebih baik, sebab tabrakan pada sudut-sudut tajam gambar sering
        terasa mengganggu pemain.
  \item \textbf{Benda yang diam cukup memerlukan \istilah{Collider}}, tanpa
        \istilah{Rigidbody}. Lantai, dinding, dan bongkahan tidak bergerak,
        sehingga tidak perlu dihitung fisika.
  \item \textbf{Benda yang bergerak memerlukan keduanya.} Tanpa
        \istilah{Rigidbody}, benda yang diberi \istilah{Collider} tetap diam
        walaupun ditabrak, sebab mesin tidak menganggapnya sebagai benda yang
        mengikuti fisika.
\end{enumerate}

\subsection{Trigger: Batas yang Tidak Menghalangi}

Sebuah \istilah{Collider} dapat ditandai sebagai \istilah{trigger}, yaitu
batas yang \textbf{tidak menghalangi} gerak benda, melainkan hanya melaporkan
bahwa ada benda lain yang memasukinya.

Perbedaan keduanya perlu dipahami dengan tepat.

\begin{itemize}
  \item \textbf{\istilah{Collider} biasa} menghalangi gerak, dan tabrakannya
        diselesaikan oleh mesin. Dipakai untuk lantai, dinding, dan rintangan.
  \item \textbf{\istilah{Trigger}} tidak menghalangi apa pun, dan hanya
        memanggil fungsi tertentu pada kode Anda. Dipakai untuk daerah yang
        menandai sesuatu: pintu yang membuka, posisi pemeriksaan
        (\istilah{checkpoint}), atau daerah yang mengakhiri permainan.
\end{itemize}

\subsection{Tiga Tipe Rigidbody}

Komponen \istilah{Rigidbody} menentukan bagaimana sebuah benda mengikuti
fisika. Tiga tipe yang tersedia diperlihatkan pada
Tabel~\ref{tab:tipe-rigidbody}.

\begin{table}[htbp]
\centering
\caption{Tiga tipe Rigidbody dan keperluannya}
\label{tab:tipe-rigidbody}
{\setlength{\tabcolsep}{4pt}\footnotesize%
\begin{tabularx}{\textwidth}{@{}L{2.6cm} L{3.6cm} Y@{}}
\toprule
\headrow \textbf{Tipe} & \textbf{Perilaku} & \textbf{Dipakai untuk} \\
\midrule
\istilah{Dynamic} & Bergerak menurut fisika: gravitasi, tumbukan, dan gaya
bekerja padanya. &
Pemain, benda yang dapat didorong, dan peluru. Inilah tipe yang paling sering
dipakai. \\
\addlinespace
\istilah{Kinematic} & Bergerak karena kode Anda, tetapi tetap dapat menabrak
benda lain. Tidak terpengaruh gravitasi. &
Benda yang geraknya Anda tentukan sendiri, misalnya pintu yang membuka atau
\istilah{platform} yang bergerak. \\
\addlinespace
\istilah{Static} & Tidak bergerak sama sekali, dan tidak memerlukan
\istilah{Rigidbody}. &
Lantai, dinding, dan seluruh benda diam. Benda bertipe ini paling murah
dihitung. \\
\bottomrule
\end{tabularx}}
\end{table}

\begin{catatabox}
Benda \istilah{static} tidak perlu diberi
\istilah{Rigidbody} sama sekali. Memberikannya justru merugikan, sebab mesin
akan menghitungnya tanpa manfaat. Kebiasaan yang benar: berikan
\istilah{Rigidbody} hanya kepada benda yang memang bergerak.
\end{catatabox}

\subsection{Pengaturan yang Sering Diubah}

Beberapa pengaturan pada \istilah{Rigidbody} perlu Anda kenal sejak awal.

\begin{itemize}
  \item \textbf{\istilah{Gravity Scale}.} Pengali percepatan gravitasi bagi
        benda itu. Nilai yang lebih besar membuat benda jatuh lebih cepat.
        Inilah padanan nilai gravitasi yang Anda tetapkan sendiri pada Bab~VI.
  \item \textbf{\istilah{Collision Detection}.} Cara mesin memeriksa tabrakan.
        Nilai \istilah{Discrete} adalah bawaannya dan paling murah, sedangkan
        \istilah{Continuous} diperlukan bagi benda yang bergerak cepat, agar
        tidak menembus benda lain.
  \item \textbf{\istilah{Interpolate}.} Menghaluskan gerak benda di antara dua
        langkah fisika. Nilainya perlu diaktifkan pada benda yang diikuti
        kamera, sebagaimana dijelaskan pada Bagian~\ref{sec:masalah-fisika}.
  \item \textbf{\istilah{Constraints}.} Membatasi gerak atau putaran pada
        sumbu tertentu. Permainan dua dimensi biasanya membekukan putaran dan
        gerak pada sumbu tegak lurus bidang.
\end{itemize}
